% Minimal retained preamble for lane-p2 cards taken from
% sources/logic-fl02-conservation-of-ramsey-s-theorem-f/pinned-original-source.tex
% (Conservation of Ramsey's theorem for pairs, arXiv:2402.11616v3).
%
% Copied verbatim from the author's own file:
%   line 2      : \documentclass{article}
%   lines 16-25 : the author's theorem declarations
%   lines 46-88 : the author's macro definitions for the reverse-mathematics principles
%   line 91     : \card
%   lines 104-107: the natural sums and products, which rely on mathabx (source line 103)
%   line 116    : \qt
%
% Omitted author preamble lines and why (nothing was renamed or redefined):
%   line 3 (graphicx), 7 (url), 9 (cancel): no figure, url or cancel use in the retained spans;
%   line 6 (cleveref): the assembler adds its own cleveref machinery, a second load clashes;
%   line 8 (\usepackage[utf8x]{inputenc}): inputenc with the utf8x option fights the forced
%       xelatex engine;
%   lines 28-37 (the rep@theorem machinery and \newreptheorem{maintheorem}): the retained spans
%       use no rep-theorem environment;
%   lines 40-43 (xcolor and the three author comment macros): not used by the retained spans;
%   lines 93-99 (the LS2 font encoding block for \smashtimes based on stix): the host TeX Live has
%       no stix directory and the author himself never uses the result -- line 105 maps \ntimes to
%       \dottimes, and line 110 shows \smashtimes already commented out in the source.
% mathabx is retained because \dotplus, \dottimes and \dotdiv come from it.
\documentclass{article}
\usepackage{amsmath,amssymb,amsthm}
\usepackage{mathabx}
\usepackage{cancel}
\usepackage{cleveref}

\newtheorem{theorem}{Theorem}
\numberwithin{theorem}{section}
\newtheorem{maintheorem}[theorem]{Main Theorem}
\newtheorem{lemma}[theorem]{Lemma}
\newtheorem{question}[theorem]{Question}
\newtheorem{proposition}[theorem]{Proposition}
\newtheorem{statement}[theorem]{Statement}
\newtheorem{remark}[theorem]{Remark}
\newtheorem{definition}[theorem]{Definition}
\newtheorem{corollary}[theorem]{Corollary}

\newcommand{\RCA}{\mathsf{RCA}}
\newcommand{\WKL}{\mathsf{WKL}}
\newcommand{\SRT}{\mathsf{SRT}}
\newcommand{\EM}{\mathsf{EM}}
\newcommand{\SEM}{\mathsf{SEM}}
\newcommand{\ADS}{\mathsf{ADS}}
\newcommand{\SCAC}{\mathsf{SCAC}}
\newcommand{\SADS}{\mathsf{SADS}}
\newcommand{\COH}{\mathsf{COH}}
\newcommand{\ACA}{\mathsf{ACA}}
\newcommand{\qvdash}{\operatorname{{?}{\vdash}}}
\newcommand{\nqvdash}{\operatorname{{?}{\nvdash}}}
\newcommand{\BSig}{\mathsf{B}\Sigma^0}
\newcommand{\ISig}{\mathsf{I}\Sigma^0}
\newcommand{\LSig}{\mathsf{L}\Sigma^0}
\newcommand{\BPi}{\mathsf{B}\Pi^0}
\newcommand{\IPi}{\mathsf{I}\Pi^0}
\newcommand{\LPi}{\mathsf{L}\Pi^0}
\newcommand{\IDel}{\mathsf{I}\Delta^0}
\newcommand{\BCDe}{\mathsf{BC}\Delta^0}
\newcommand{\WF}{\mathsf{WF}}
\newcommand{\RT}{\mathsf{RT}}
\newcommand{\BME}{\mathsf{BME}}
\newcommand{\TT}{\mathsf{TT}}
\renewcommand{\O}{\mathcal{O}}
\renewcommand{\L}{\mathcal{L}}
\newcommand{\C}{\mathcal{C}}
\newcommand{\R}{\mathcal{R}}
\renewcommand{\S}{\mathcal{S}}
\newcommand{\T}{\mathcal{T}}
\renewcommand{\P}{\mathcal{P}}
\newcommand{\M}{\mathcal{M}}
\newcommand{\Nc}{\mathcal{N}}
\renewcommand{\P}{\mathcal{P}}
\newcommand{\Q}{\mathcal{Q}}
\newcommand{\NN}{\mathbb{N}}
\newcommand{\PP}{\mathbb{P}}
\newcommand{\uh}[0]{{\upharpoonright}}
\newcommand{\card}{\operatorname{card}}
\newcommand{\nplus}{\dotplus}
\newcommand{\ntimes}{\dottimes}
\newcommand{\nminus}{\dotdiv}
\newcommand{\nsum}{\dot{\sum}}
\def\qt#1{``#1''}
\title{The forcing question for split pairs is decided Y-prime computably}
\author{Quentin Le Houérou, Ludovic Levy Patey, Keita Yokoyama}
\date{arxiv:2402.11616v3; CC BY 4.0}
\begin{document}
\maketitle
\noindent\textit{Editorial note.} The counted unit is the article's own Lemma lem:sads-forcing-question-sigma01 at source lines 1149-1156 together with the proof environment that begins two lines after the statement ends (1158-1166); the pairing is the ordinary adjacent one.\par\medskip
\noindent\textit{Editorial note.} All five context spans lie inside the proof environment of the outer Proposition prop:sads-wf, which opens at source line 1111: they are the definitions of the sets U and V, of a condition, of a split pair, of the two forcing relations and of the forcing question. They are retained as context because the counted lemma is unreadable without them, and context spans are not counted units; the outer proposition itself is not counted by this lane, because its interval is nested with the counted lemma and nested spans are counted once.\par\medskip
\noindent\textit{Editorial note.} The author records in a footnote at source line 1141 that the proof of lem:sads-wf-finite-tree-forcing was flawed in the published version of this article at Transactions of the American Mathematical Society 378 (2025) no. 3, pp. 2157-2186, and that he fixed it here by introducing the forcing question and proving this lemma. The footnote is not part of the retained material but the fact is recorded here for the reader.\par\medskip
\noindent\textit{Editorial note.} The retained preamble reproduces verbatim the author's document class at source line 2, his theorem declarations at source lines 16-25, his macro definitions at source lines 46-88 and 91, and his natural-operation macros at source lines 104-107, and it retains mathabx and cleveref; the stix-dependent LS2 block and the inputenc load with the utf8x option are not retained.\par\medskip
\noindent\textit{Editorial note.} No retained line contains a citation, so the card carries no bibliography. The article's conservation theorems and its other lemmas are not counted here.\par\medskip
\noindent\textit{Allowed standard prerequisites (editorial).}\begin{itemize}
\item Conditions, extensions and split pairs of the SADS forcing. For a linear order of order type omega plus omega* and a model with no solution, the elements split into the set U of those with finitely many predecessors and the set V of those with finitely many successors; a condition is a pair of a finite ascending sequence inside U and a finite descending sequence inside V, ordered by extension of both components, and a split pair is a pair whose ascending part lies entirely below its descending part; the author's own account is retained at source lines 1112, 1114 and 1117-1119.. Fixes all the objects of the counted lemma, and in particular the observation that a split pair has one component inside U or inside V, which is what makes the first case of the proof work.
\item The two forcing relations and the forcing question. A condition forces convergence on side i when the corresponding finite string already converges, and forces divergence on side i when no compatible split pair extension converges there; the forcing question between the two sides holds when some split pair extending the condition makes both designated computations converge; the author's own definitions are retained at source lines 1121-1127 and 1143-1145.. The counted lemma is exactly the statement that this forcing question is decided on a non-empty condition, with the two cases producing extensions forcing convergence or divergence on one side.
\end{itemize}\medskip
Let $L = (M, \leq_L)$ be such an order. If $\M$ already contains an $M$-unbounded ascending or descending sequence for $\leq_L$, then we are done, so assume otherwise. Let $U \subseteq M$ and $V \subseteq M$ be the sets of elements that have a finite amount of predecessors and successors, respectively. By assumption, $U \sqcup V = M$, and $U$ and $V$ are both $\Delta^0_2(Y)$.

A condition is a pair $(\sigma, \tau)$ with $\sigma$ an $M$-finite $L$-ascending sequence included in $U$ and $\tau$ an $M$-finite $L$-descending sequence included in $V$. We have $(\sigma_2, \tau_2) \leq (\sigma_1, \tau_1)$ if $\sigma_2 \succeq \sigma_1$ and $\tau_2 \succeq \tau_1$.

\begin{definition}
A \emph{split pair} is a pair $(\sigma, \tau)$ of strings such that $\sigma$ is an $M$-finite $L$-ascending sequence, $\tau$ is an $M$-finite $L$-descending sequence, and $\max_L \sigma <_L \min_L \tau$.
\end{definition}

\begin{definition}
Let $(\sigma_0, \sigma_1)$ be a condition, $e, x \in M$ and $i < 2$, we write:
\begin{itemize}
    \item $(\sigma_0, \sigma_1) \Vdash_i \Phi_e^{G \oplus Y}(x) \downarrow$ if $\Phi_e^{\sigma_i \oplus Y}(x) \downarrow$. 
    \item $(\sigma_0, \sigma_1) \Vdash_i \Phi_e^{G \oplus Y}(x) \uparrow$ if $\forall \sigma'_i \succeq \sigma_i$ such that $(\sigma'_i, \sigma_{1-i})$ is a split pair, $\Phi_e^{\sigma'_i \oplus Y}(x)\uparrow$.
\end{itemize}   
\end{definition}

\begin{definition}
    Let $(\sigma, \tau)$ be a condition and $e_0, e_1, x_0, x_1 \in M$. We write $(\sigma, \tau) \qvdash \Phi_{e_0}^{G_0 \oplus Y}(x_0) \downarrow \vee \Phi_{e_1}^{G_1 \oplus Y}(x_1) \downarrow$ if there exists some split pair $(\sigma', \tau')$ such that $\sigma \preceq \sigma'$, $\tau \preceq \tau'$, $\Phi_{e_0}^{\sigma' \oplus Y}(x_0) \downarrow$ and $\Phi_{e_1}^{\tau' \oplus Y}(x_1) \downarrow$. 
\end{definition}

\begin{lemma}\label[lemma]{lem:sads-forcing-question-sigma01}
    Let $(\sigma_0, \tau_0)$ be a condition, and let $e_0, e_1, x_0, x_1 \in M$. Then:
    \begin{itemize}
        \item If $(\sigma_0, \tau_0) \qvdash \Phi_{e_0}^{G_0 \oplus Y}(x_0) \downarrow \vee \Phi_{e_1}^{G_1 \oplus Y}(x_1) \downarrow$, then there exist an extension $(\sigma_1, \tau_1) \leq (\sigma_0, \tau_0)$ and some $i < 2$ such that $(\sigma_1, \tau_1) \Vdash_i \Phi_{e_i}^{G_i \oplus Y}(x_i) \downarrow$.
        \item If $(\sigma_0, \tau_0) \nqvdash \Phi_{e_0}^{G_0 \oplus Y}(x_0) \downarrow \vee \Phi_{e_1}^{G_1 \oplus Y}(x_1) \downarrow$, then there exist an extension $(\sigma_1, \tau_1) \leq (\sigma_0, \tau_0)$ and some $i < 2$ such that $(\sigma_1, \tau_1) \Vdash_i \Phi_{e_i}^{G_i \oplus Y}(x_i) \uparrow$.
    \end{itemize}
    Furthermore, deciding which case holds and computing the extension $(\sigma_1, \tau_1)$ can be done $Y'$-computably.
\end{lemma}

\begin{proof}
Assume first that $(\sigma_0, \tau_0) \qvdash \Phi_{e_0}^{G_0 \oplus Y}(x_0) \downarrow \vee \Phi_{e_1}^{G_1 \oplus Y}(x_1) \downarrow$. Then there exists a split pair $(\sigma_2, \tau_2)$ such that $\sigma_0 \preceq \sigma_2$, $\tau_0 \preceq \tau_2$, $\Phi_{e_0}^{\sigma_2 \oplus Y}(x_0)\downarrow$ and $\Phi_{e_1}^{\tau_2 \oplus Y}(x_1)\downarrow$. Using $Y$, we can find such a split pair. Since $(\sigma_2, \tau_2)$ is a split pair, either $\sigma_2 \subseteq U$ or $\tau_2 \subseteq V$, and deciding which case holds can be done $Y'$-computably. If $\sigma_2 \subseteq U$, then letting $\sigma_1 := \sigma_2$ and $\tau_1 = \tau_0$ yields an extension $(\sigma_1, \tau_1) \leq (\sigma_0, \tau_0)$ such that $(\sigma_1, \tau_1) \Vdash_0 \Phi_{e_0}^{G_0 \oplus Y}(x_0) \downarrow$. The case $\tau_2 \subseteq V$ is symmetrical.  \\

Now suppose that $(\sigma_ 0, \tau_ 0) \nqvdash \Phi_{e_0}^{G_0 \oplus Y}(x_0) \downarrow \vee \Phi_{e_1}^{G_1 \oplus Y}(x_1) \downarrow$. Consider the program which runs through all $M$-finite $L$-ascending sequences $\sigma_2$ such that $\sigma_ 0 \preceq \sigma_2$, $\max_L \sigma_2 <_L \min_L \tau_0$ and $\Phi_{e_0}^{\sigma_2 \oplus Y}(x_0) \downarrow$. Whenever a new $\sigma_2$ is found whose value of $\max_L \sigma_2$ is smaller than all previously found ones, enumerate $\max_L \sigma_2$. By construction, this program yields a $Y$-computable $L$-descending sequence. By our assumption, $\M$ does not contain any solution to the given instance of $\SADS$, this enumeration must therefore be finite. Hence, either no such $\sigma_2$ is enumerated, in which case $(\sigma_ 0, \tau_ 0) \Vdash_0 \Phi_{e_0}^{G_0 \oplus Y}(x_0) \uparrow$ and we are done, or else we can $Y'$-computably find some $\sigma_2$ such that $\sigma_0 \preceq \sigma_2$, $\max_L \sigma_2 <_L \min_L \tau_0$ and $\Phi_{e_0}^{\sigma_2 \oplus Y}(x_0) \downarrow$ and for which $\max_L \sigma_2$ is minimal. 

Similarly, either we already have $(\sigma_ 0, \tau_ 0) \Vdash_1 \Phi_{e_1}^{G_1 \oplus Y}(x_1) \uparrow$ or we can $Y'$-computably find some $M$-finite $L$-descending sequence $\tau_2$ such that $\tau_0 \preceq \tau_2$, $\min_L \tau_2 >_L \max_L \sigma_0$ and $\Phi_{e_1}^{\tau_2 \oplus Y}(x_1) \downarrow$ and such that $\min_L \tau_2$ is maximal among all such sequences. 

In the case where such $\sigma_2$ and $\tau_2$ are found, we cannot have $\sigma_2 \subseteq U$ and $\tau_2 \subseteq V$, since then $(\sigma_2, \tau_2)$ would form a split pair, contradicting our assumption that $(\sigma_0, \tau_0) \nqvdash \Phi_{e_0}^{G_0 \oplus Y}(x_0) \downarrow \vee \Phi_{e_1}^{G_1 \oplus Y}(x_1) \downarrow$, and we can $Y'$-computably find whether $\sigma_2 \not \subseteq U$ or $\tau_2 \not \subseteq V$. Without any loss of generality, assume $\sigma_2 \not \subseteq U$. Define $\sigma_1 = \sigma_0$ and $\tau_1 = \tau_0 \cup \{\max_L \sigma_2\}$, then $(\sigma_1, \tau_1)$ is a condition extending $(\sigma_0, \tau_0)$. Indeed, since $\sigma_2 \not \subseteq U$, we have $\max_L \sigma_2 \in V$, and by construction $\max_L \sigma_2 <_L \min_L \tau_0$. Finally, we have $(\sigma_1, \tau_1) \Vdash_0 \Phi_{e_0}^{G_0 \oplus Y}(x_0) \uparrow$. Indeed, by minimality of $\max_L \sigma_2$, there is no $M$-finite $L$-ascending sequence $\sigma_2'$ such that $\sigma_1 = \sigma_0 \preceq \sigma_2'$, $\max_L \sigma_2' <_L \min_L \tau_1 = \max_L \sigma_2$ and $\Phi_{e_0}^{\sigma_2' \oplus Y}(x_0) \downarrow$.
\end{proof}

\end{document}
